\lstset{style=Python,escapeinside={(*@}{@*)}}
\fboxsep=.6pt

% Define a toggle for exerciseonly mode (false by default)
\newtoggle{exerciseonly}
\settoggle{exerciseonly}{false}

% Check for the class option 'exerciseonly'
\makeatletter
\@ifclasswith{book}{exerciseonly}{%
	\settoggle{exerciseonly}{true}%
}{}
\makeatother

% load package if document contains at least one citation
\usepackage[maxbibnames=99,sorting=none,style=numeric,backend=biber]{biblatex}
% enable option to print bibliography as a numbered subsection
\defbibheading{bibsubsection}[\bibname]{%
  \subsection{#1}%
}
\addbibresource{refs.bib}

\begin{document}
%% Project-specific Preambles 
\settoggle{includelogo}{true}

\newcommand{\thesubject}{Informatik}
\newcommand{\thetopic}{Programmieren}
\newcommand{\theplace}{Winterthur}
\newcommand{\thelogo}{Logos/FreeFlower.pdf}
\newcommand{\logosize}{.5\textwidth}

\settoggle{emoji}{true}

% macro for creating a title page
\renewcommand{\maketitlepage}[2]{
	% arg1: title
	% arg2: subtitle
	\begin{titlepage}
		\vspace*{0.25\textheight}
		\begin{center}
			% logo
			\iftoggle{includelogo}{\begin{figure}[H]\centering{\includegraphics[width=7cm]{Logos/Institutions/LeeLogo.pdf}}
				\end{figure}}{}
			\thesubject\\
			\begin{tcolorbox}[enhanced,colframe=RoyalBlue!80,colback=RoyalBlue!20, hbox,lifted shadow={1mm}{-2mm}{3mm}{0.1mm}%
						{Black!50!white}]
				\begin{varwidth}{\textwidth}
					\begin{center}
						\ifblank{#2}{
							% if no subtitle is given (#2 blank)
							{\huge \textbf{#1}}\\
							\iftoggle{emoji}{\makeSpiral}{}
							{\large \textbf{Skript}}\\[0.2cm]
						}{
							% if #2 not blank
							{\huge \textbf{#1}}\\
							\iftoggle{emoji}{\makeSpiral}{}
							{\large \textbf{#2}}\\[0.2cm]
						}
					\end{center}
				\end{varwidth}
			\end{tcolorbox}


			{\large Kantonsschule im Lee, \today}

			%      \tikz[remember picture,overlay] \node[opacity=0.3,inner sep=0pt] at (current page.center){\includegraphics[width=\paperwidth,height=\paperheight]{LeeLogo}};

		\end{center}
	\end{titlepage}
}
\lhead[ \leftmark ]{\thetopic}
\rhead[\thesubject]{\textcopyleft{} \thesubject, \the\year}

\clearpage
\def\arrayB{"grinning-face","bar-chart","laptop","technologist","robot","woman-student","hundred-points"}

% special case to make "diplomas" using the title page template
\foreach \person/\platz/\mention in {
		% Livio Adorni/Honorable Mention/,
		% Jonas Billeter/Honorable Mention/,
		Lars Bärlocher/3. Platz/{herausragendes }}{
		\maketitlepage{\textbf{\platz\\[.3cm]\person}}{Informatik-Interessenwoche zum Thema Web-Technologien\\[.3cm]Für ein \mention Projekt mit HTML, CSS und JS}
	}
\end{document}